\chapter{Floquet theory for periodically driven systems}\label{ch:floquet}

This chapter introduces the physics of periodically driven quantum systems and
develops Floquet theory from its basic principles. Extensive reviews are available in
Refs.~\cite{grifoni1998driven,kohler2005driven,chu2004beyond,holthaus2016floquet,eckardt2017colloquium,oka2019floquet};
here we develop the concepts needed to understand and use the Floquet methods in \pkg{}.

\section{Periodic Hamiltonians and the Floquet theorem}

Consider a Hamiltonian with period $T=2\pi/\Omega$,
\begin{equation}
  H(t+T)=H(t),\qquad \ii\partial_t\ket{\psi(t)}=H(t)\ket{\psi(t)} .
\end{equation}
The system exchanges energy with the drive, and its state evolves in time. The
periodicity constrains the time-evolution operator $U(t,t_0)$: the evolution
over one period is identical for starting times separated by an integer number of periods.
The one-period propagator $U(T,0)$ is unitary; let its eigenvectors be $\ket{\varphi_\alpha}$
with eigenvalues $\ee^{-\ii\qe_\alpha T}$. The states
\begin{equation}\label{eq:floquet-state}
  \ket{\psi_\alpha(t)}=U(t,0)\ket{\varphi_\alpha}=\ee^{-\ii\qe_\alpha t}\ket{u_\alpha(t)},
  \qquad \ket{u_\alpha(t+T)}=\ket{u_\alpha(t)} ,
\end{equation}
are solutions of the Schr\"odinger equation that differ from stationary states only by a
periodic \enquote{micromotion} $\ket{u_\alpha(t)}$. This is the Floquet
theorem~\cite{floquet1883equations,shirley1965solution,zeldovich1967quasienergy,ritus1967shift}.
The quantities $\qe_\alpha$ are the \emph{quasienergies}, the $\ket{u_\alpha(t)}$ are
the \emph{Floquet modes}, and the $\ket{\psi_\alpha(t)}$ are the \emph{Floquet states}. The
modes form a complete basis at every time, so any solution is
$\ket{\psi(t)}=\sum_\alpha c_\alpha\ee^{-\ii\qe_\alpha t}\ket{u_\alpha(t)}$ with constant
coefficients: the Floquet states play the role of the energy eigenstates of a static
problem.

The quasienergy is defined only modulo $\Omega$. Replacing $\qe_\alpha$ by
$\qe_\alpha+\ell\Omega$ and $\ket{u_\alpha(t)}$ by $\ee^{\ii\ell\Omega t}\ket{u_\alpha(t)}$
leaves $\ket{\psi_\alpha(t)}$ unchanged for every integer $\ell$. This is the temporal
analogue of the quasimomentum in Bloch's theorem: quasienergies can be represented in a Brillouin zone of
width $\Omega$, for instance $[-\Omega/2,\Omega/2)$.

\begin{physics}
A static system has energy levels; a periodically driven system has quasienergy levels.
They determine the phases accumulated over each period and enter the conditions for
interference, resonances and spectroscopic transitions. A state that is
an energy eigenstate without the drive is continuously deformed into a Floquet state as the
drive is switched on, as long as no resonance is crossed.
\end{physics}

\paragraph{Effective Hamiltonian and micromotion.}
The one-period propagator can be written as $U(T,0)=\ee^{-\ii H_FT}$ with the
\emph{Floquet Hamiltonian} $H_F$, whose eigenvalues are the quasienergies. $H_F$ depends on
the time at which the period starts, while its spectrum is time-origin invariant. The full evolution
separates into a slow evolution under a time-independent \emph{effective Hamiltonian}
$\Heff$ and a periodic micromotion generated by the kick operator
$\mathcal K(t)=\mathcal K(t+T)$,
$U(t,0)=\ee^{-\ii\mathcal K(t)}\ee^{-\ii\Heff t}\ee^{\ii\mathcal K(0)}$. At $t=T$ this gives
$H_F=\ee^{-\ii\mathcal K(0)}\Heff\ee^{\ii\mathcal K(0)}$: both operators have the same
spectrum, and $\Heff$ is invariant under shifts of the time
origin~\cite{goldman2014periodically,bukov2015universal,eckardt2015high}. The
high-frequency expansion of \cref{sec:van-vleck} approximates $\Heff$.

\section{Fourier space: the Sambe picture}\label{sec:sambe}

Since $H(t)$ and $\ket{u(t)}$ are periodic, both can be expanded in harmonics of $\Omega$.
\pkg{} uses the convention
\begin{equation}\label{eq:harmonics}
  H(t)=\sum_k H_k\,\ee^{\ii k\Omega t},\qquad H_{-k}=H_k\dg,\qquad
  \ket{u(t)}=\sum_m \ee^{\ii m\Omega t}\ket{u_m}.
\end{equation}
Inserting $\ket{\psi(t)}=\ee^{-\ii\qe t}\ket{u(t)}$ into the Schr\"odinger equation and
comparing the coefficients of $\ee^{\ii m\Omega t}$ gives
\begin{equation}\label{eq:sambe}
  \sum_{m'}\left(H_{m-m'}+m\Omega\,\delta_{mm'}\right)\ket{u_{m'}}=\qe\,\ket{u_m}.
\end{equation}
This is a time-independent eigenvalue problem for the infinite block matrix
\begin{equation}\label{eq:sambe-matrix}
  Q=\begin{pmatrix}
    \ddots & & & & \\
    & H_0-\Omega & H_{-1} & H_{-2} & \\
    & H_1 & H_0 & H_{-1} & \\
    & H_2 & H_1 & H_0+\Omega & \\
    & & & & \ddots
  \end{pmatrix},
\end{equation}
which acts in the extended (Sambe) space of $T$-periodic states with scalar product
$\langle\!\langle a\vert b\rangle\!\rangle=T^{-1}\int_0^T\dd t\,\braket{a(t)}{b(t)}$~\cite{sambe1973steady}.
The diagonal blocks are copies of the time-averaged Hamiltonian $H_0$ shifted by multiples
of $\Omega$; the off-diagonal blocks are the drive harmonics. For a monochromatic drive
linear in a parameter, only $H_{\pm1}$ are nonzero and $Q$ is block-tridiagonal. \pkg{}
truncates the matrix to the blocks $\abs m\le M$, where $M$ is the \emph{Floquet cutoff}
(option \code{floquet_cutoff}).

\paragraph{Replicas.}
The eigenvalue problem \eqref{eq:sambe} is invariant under the shift
$\ket{u_m}\to\ket{u_{m-\ell}}$ combined with $\qe\to\qe+\ell\Omega$, for any integer $\ell$.
Every physical Floquet state therefore appears as an infinite ladder of \emph{replicas},
one per Brillouin zone. The harmonic weights $\norm{u_m}^2$ sum to one (for
$\langle\!\langle u\vert u\rangle\!\rangle=1$), and their centroid
$\expval m=\sum_m m\norm{u_m}^2$ shifts by $\ell$ from one replica to the next. The
truncation distorts the replicas near the edges of the window $\abs m\le M$. The code
therefore works with the most central replica and reports the weight in the outermost
blocks as a convergence diagnostic.

\paragraph{Mean energy.}
Taking the expectation value of \eqref{eq:sambe} in a normalized Floquet mode gives
\begin{equation}\label{eq:mean-energy}
  \bar E\equiv\frac1T\int_0^T\dd t\,\matel{\psi(t)}{H(t)}{\psi(t)}=\qe-\Omega\expval m ,
\end{equation}
the time-averaged energy of the Floquet state. The value of $\bar E$ is replica
invariant, making it a natural label for a Floquet state; it is reported as
\code{mean_energies}.

\section{An exactly solvable example}\label{sec:circular}

Consider a spin $\tfrac12$ (or any two-level system) in a static field and a circularly
polarized field,
\begin{equation}
  H(t)=\frac{\omega_0}2\sigma_z+\frac b2\left(\ee^{-\ii\Omega t}\sigma_++\ee^{\ii\Omega t}\sigma_-\right),
  \qquad H_{-1}=\frac b2\sigma_+,\quad H_1=\frac b2\sigma_-,
\end{equation}
This problem is solved by transforming to a rotating frame. The transformation
$W(t)=\ee^{\ii\Omega t\sigma_z/2}$ gives, by \cref{eq:gauge}, the static Hamiltonian
$H'=\tfrac12(\omega_0-\Omega)\sigma_z+\tfrac b2\sigma_x$ with eigenvalues $\pm\tfrac12w$,
$w=\sqrt{(\omega_0-\Omega)^2+b^2}$. Transforming back, the spin-up component of each solution
acquires $\ee^{-\ii\Omega t/2}$ and the spin-down component $\ee^{+\ii\Omega
t/2}=\ee^{-\ii\Omega t/2}\ee^{\ii\Omega t}$. The Floquet states therefore have
\begin{equation}\label{eq:circular}
  \qe_\pm=\frac{\Omega\pm w}2\pmod\Omega,\qquad
  \ket{u(t)}=\ket{u_0}+\ee^{\ii\Omega t}\ket{u_1},\quad \ket{u_0}\propto\ket\up,\ \ket{u_1}\propto\ket\dn .
\end{equation}
Exactly two harmonics have nonzero amplitudes. In the photon picture of \cref{sec:photons}, the
spin-down component is accompanied by one more photon in the field than the spin-up
component. Away from resonance, $\abs{\omega_0-\Omega}\gg b$, the two quasienergies
approach $\omega_0/2$ and $-\omega_0/2+\Omega$, which are the bare energy of $\ket\up$ and the
bare energy of $\ket\dn$ plus one photon. At resonance, $\Omega=\omega_0$, these replicas
would cross; the drive opens a gap $b$ between them, the Rabi splitting. A system prepared in
$\ket\up$ oscillates as $P_\dn(t)=(b^2/w^2)\sin^2(wt/2)$, so the Rabi splitting at resonance
equals the frequency of the Rabi oscillations. For a linearly polarized field,
$b\cos(\Omega t)\sigma_x$, the same picture holds up to counter-rotating corrections of
order $b^2/\omega_0$ (\cref{ch:ex-tls}). \pkg{} reproduces \cref{eq:circular} to machine
precision with all its methods, and the code is tested against it.

\section{The photon picture}\label{sec:photons}

Equation~\eqref{eq:sambe} has a transparent interpretation. Consider the system coupled to
a single quantized field mode of frequency $\Omega$, $H_{\mathrm{sys}}+\Omega a\dg
a+H_{\mathrm{int}}$, in a coherent state with a large mean photon number $\bar N$. In the
limit $\bar N\to\infty$ at fixed field amplitude, the Hamiltonian in the basis
$\ket{\alpha}\otimes\ket{\bar N+m}$ becomes exactly $Q$. Here $m$ is the change of the photon
number relative to $\bar N$: $\abs m$ photons have been absorbed from the field for $m<0$
and emitted into it for
$m>0$~\cite{shirley1965solution,cohentannoudji1977dressed,cohentannoudji1998atom}. The Floquet
states are the \emph{dressed states} of the system and the field. The diagonal blocks
$H_0+m\Omega$ are the bare levels with $m$ extra photons, and the harmonic weight
$\norm{u_m}^2$ is the probability that $m$ photons have been exchanged with the field. The
quasienergy spectrum is the spectrum of the dressed system, measured relative to a field
energy that is defined only up to multiples of $\Omega$.

\begin{physics}
Two bare levels $E_a$ and $E_b$ are resonant when $E_b-E_a=k\Omega$ for some integer $k$:
in the photon picture, the states $\ket a\otimes\ket{\bar N}$ and $\ket b\otimes\ket{\bar
N-k}$ are degenerate, and absorbing $k$ photons transfers the system from $a$ to $b$ without
energy cost. The drive couples them, and the Floquet spectrum shows an avoided crossing, the
Autler--Townes splitting~\cite{autler1955stark}. For $k=1$ its size is set directly by the
first harmonic of the drive. For $k>1$ it comes from the harmonic $H_k$, from virtual
intermediate states, or, for strong modulation, from Bessel functions of the modulation
index.
\end{physics}

\section{Gauge freedom and the energy zero}\label{sec:gauge}

A time-dependent unitary transformation $\ket{\psi}\to W(t)\ket\psi$ maps the Schr\"odinger
equation onto one with
\begin{equation}\label{eq:gauge}
  H'(t)=W(t)H(t)W\dg(t)+\ii\,\bigl(\partial_tW(t)\bigr)W\dg(t).
\end{equation}
If $W(t)$ is periodic, the quasienergies of $H$ and $H'$ coincide modulo $\Omega$, while the
Floquet modes are multiplied by $W(t)$. Two consequences are used throughout this manual.
\begin{enumerate}
\item A c-number term $c(t)$ in $H(t)$ multiplies $\psi$ by the phase
$\exp[-\ii\int c\,\dd t]$. Its time average shifts all quasienergies equally; its
oscillating part only redistributes the harmonic weights. \pkg{} keeps such terms by
default, so the Floquet problem uses $H(t)$ exactly as built (option
\code{drop_scalar_harmonics} removes the oscillating c-number terms).
\item A time-dependent superconducting phase on the tunnelling terms can be traded for a
time-dependent electrochemical potential, i.e.\ a voltage (\cref{sec:phase-voltage}).
\end{enumerate}

\section{Regimes}\label{sec:regimes}

\paragraph{Weak driving.}
Without the drive, $\qe_a=E_a$ and $\ket{u_m}=\delta_{m0}\ket a$. To second order in the
drive, the quasienergy of a nondegenerate state shifts by the dynamical (AC) Stark shift
\begin{equation}\label{eq:ac-stark}
  \delta\qe_a=\sum_{(b,m)\ne(a,0)}\frac{\abs{\matel b{H_m}a}^2}{E_a-E_b-m\Omega},
\end{equation}
obtained by perturbation theory in the Sambe space. Here $E_a$ and $\ket a$ are the
eigenvalues and eigenstates of the time-averaged Hamiltonian $H_0$. For a drive linear in a
parameter, $H_0$ is the undriven Hamiltonian. For a nonlinear drive, such as the phase drive,
$H_0$ itself changes at second order in the amplitude (\cref{eq:phase-h0}), and this change
adds to \eqref{eq:ac-stark}. The denominators vanish at the resonances
$E_b-E_a=-m\Omega$, where degenerate perturbation theory, the rotating-wave approximation
of \cref{sec:rwa}, must be used.

\paragraph{Resonant driving.}
Near a one-photon resonance the two coupled states $\ket a\otimes\ket{\bar N}$ and $\ket
b\otimes\ket{\bar N-1}$ form an isolated two-level system. Its splitting at resonance, the
Rabi splitting, is $2\abs{\matel b{H_{-1}}a}$; it equals the frequency of Rabi
oscillations between $a$ and $b$~\cite{rabi1937space}. The counter-rotating
couplings shift the resonance by an amount of second order in the drive, the
Bloch--Siegert shift~\cite{bloch1940magnetic,shirley1965solution}.

\paragraph{Slow driving.}
For $\Omega$ small compared with the gaps to all other levels, the system follows the
instantaneous eigenstate $\ket{a(t)}$ adiabatically. The quasienergy becomes the time
average of the instantaneous energy $E_a(t)$, corrected by the geometric (Berry) phase
$\gamma_a$ accumulated per period~\cite{breuer1989adiabatic,drese1999floquet},
\begin{equation}\label{eq:adiabatic}
  \qe_a\simeq\frac1T\int_0^T E_a(t)\,\dd t-\frac{\gamma_a}{T}\pmod\Omega .
\end{equation}
The Berry phase vanishes when the parameter moves back and forth along a line, as for a
single-parameter cosine drive.

\paragraph{Fast driving.}
For $\Omega$ larger than all other energy scales the time-averaged Hamiltonian governs
the leading response. The quasienergies then follow from a static effective
Hamiltonian, expanded in powers of $1/\Omega$
(\cref{sec:van-vleck})~\cite{magnus1954exponential,maricq1982application,rahav2003effective,eckardt2015high,bukov2015universal,mikami2016brillouin}.
Strong fast driving renormalizes couplings by Bessel functions, which can switch them off
entirely: coherent destruction of tunnelling~\cite{grossmann1991coherent} and dynamical
localization~\cite{dunlap1986dynamic} (\cref{ch:ex-highfreq}).

\paragraph{Strong driving.}
When the modulation of a level is large compared with $\Omega$, many harmonics are
involved and multiphoton resonances appear, with couplings renormalized by Bessel functions
of the modulation index. In two-level systems this produces the
Landau--Zener--St\"uckelberg interference
patterns~\cite{shevchenko2010landau,oliver2005mach,sillanpaa2006continuous,ivakhnenko2023nonadiabatic}
(\cref{ch:ex-tls}).

\section{Floquet states coupled to a continuum}\label{sec:continuum}

The subgap states of a junction are discrete, but the continuum begins at
$E_{\mathrm{ground}}+\Delta$. A replica $\ket a\otimes\ket{\bar N+m}$ of a bound state that
falls into the continuum, $E_a+m\Omega>E_{\mathrm{ground}}+\Delta$, is degenerate with
continuum states. In an infinite system it acquires a decay rate given by the golden rule,
$2\pi\abs{V}^2\rho_c$, with $V$ the drive matrix element and $\rho_c$ the density of
continuum states. Physically, the drive can excite quasiparticles out of the bound state,
the analogue of the multiphoton ionization of an atom. The Floquet state then becomes a
resonance with a finite lifetime~\cite{chu2004beyond,moiseyev2011non}.

In a finite bath the continuum is a set of discrete levels. The replicas of these levels
cross the bound state as a parameter is varied, and every crossing becomes a narrow avoided
crossing. As the bath is refined, the crossings become denser and narrower; their positions
and gaps depend on the discretization. Hone, Ketzmerick and Kohn analysed this generic
behaviour of Floquet spectra with dense sets of levels, and the related absence of a strict
adiabatic limit~\cite{hone1997time,hone2009statistical}.

\begin{caution}
Narrow avoided crossings with continuum replicas redistribute a dressed bound
state's weight among the discrete finite-bath states. Their positions and gaps
depend on the discretization; in the continuum their collective effect gives a
finite lifetime. Identify physical features by following the dressed states'
static weights and checking stability under bath refinement (\cref{ch:ex-numerics}).
\end{caution}

\section{What is observable}

Weak-probe spectroscopy measures quasienergy differences combined with integer
multiples of the drive frequency. The probe induces transitions between Floquet states at the
frequencies
\begin{equation}\label{eq:transitions}
  \omega=\qe_\alpha-\qe_\beta+q\Omega,\qquad q\in\mathbb Z ,
\end{equation}
with strengths given by the harmonics of the probe matrix element between the two Floquet
modes. Which lines are actually seen also depends on the populations of the Floquet
states. Relaxation sets these populations: in an open system they follow from a Floquet
master equation~\cite{blumel1991dynamical,kohler1997floquet,breuer2007theory,hone2009statistical}.
\pkg{} computes the quasienergies, modes and probe matrix elements
(\cref{ch:probes}). A calculation of the populations using the same modes is planned.
At present, spectra are evaluated for explicitly specified initial states.

\begin{incode}
\code{solve_floquet} returns \code{FloquetStates} with the \code{quasienergies} on a definite
branch (\cref{sec:sambe-method}), the \code{folded} values in a chosen zone, the
\code{mean_energies} \eqref{eq:mean-energy}, the harmonics $\ket{u_m}$ of the modes
(\code{modes}), the \code{harmonic_weights} $\norm{u_m}^2$ and the \code{static_weights},
the time-averaged weights $\sum_m\abs{\braket a{u_m}}^2$ of the undriven eigenstates
$\ket a$. \code{mode_at(t)} and \code{state_at(t)} evaluate $\ket{u(t)}$ and $\ket{\psi(t)}$.
\code{matrix_elements} computes probe matrix elements (\cref{ch:probes}), and
\code{dressed_state(label)} returns the index of the Floquet state that continues the
static state \code{label}.
\par\smallskip\usermanual{docs/floquet.md}.
\end{incode}
